\documentclass[12pt]{article}
\usepackage[a4paper, margin=1in]{geometry}
\usepackage{amsmath, amssymb}
\usepackage{enumitem}

\title{\textbf{Lecture Title: Linear Transformations and\\
Expectations of Multiple Random Variables}}
\author{Course: CSE400 - Fundamentals of Probability in Computing\\
Lecture Number: 20}
\date{}

\begin{document}
\maketitle

\section*{1. Overview}
This lecture focuses on extending probability concepts from single random variables to multiple random variables using vector representation. The key ideas include:
\begin{itemize}
\item Representation of multiple random variables as a random vector.
\item Understanding mean vector, correlation matrix, and covariance matrix.
\item Studying linear transformations of random vectors.
\item Analyzing how statistical properties such as mean and correlation transform under linear mappings.
\item Introduction to multivariate distributions and their properties.
\end{itemize}

\section*{2. Definitions and Notation}

\textbf{Random Vector}

For random variables $X_1, X_2, \dots, X_N$, we define a random vector:
\[
X = 
\begin{bmatrix}
X_1 \\
X_2 \\
\vdots \\
X_N
\end{bmatrix}
\]
$X$ is called a random vector. It provides a compact representation of multiple random variables.

\textbf{Mean Vector}

The mean vector of $X$ is defined as:
\[
\mu_X = E[X]
\]

Expanded form:
\[
\mu_X =
\begin{bmatrix}
E[X_1] \\
E[X_2] \\
\vdots \\
E[X_N]
\end{bmatrix}
\]

\textbf{Correlation Matrix}

The correlation matrix is defined as:
\[
R_{XX} = E[XX^T]
\]

Element-wise:
\[
R_{XX}(i,j) = E[X_i X_j]
\]

Matrix form:
\[
R_{XX} =
\begin{bmatrix}
E[X_1 X_1] & E[X_1 X_2] & \cdots & E[X_1 X_N] \\
E[X_2 X_1] & E[X_2 X_2] & \cdots & E[X_2 X_N] \\
\vdots & \vdots & \ddots & \vdots \\
E[X_N X_1] & E[X_N X_2] & \cdots & E[X_N X_N]
\end{bmatrix}
\]

\textbf{Covariance Matrix}

\[
C_{XX} = E[(X - \mu_X)(X - \mu_X)^T]
\]

Expanded interpretation:
\begin{itemize}
\item Diagonal elements: $\mathrm{Var}(X_i)$
\item Off-diagonal elements: $\mathrm{Cov}(X_i, X_j)$
\item Symmetry property: $\mathrm{Cov}(X_i, X_j) = \mathrm{Cov}(X_j, X_i)$
\end{itemize}

\section*{3. Assumptions and Conditions}

\begin{itemize}
\item Random variables are assumed to be well-defined with finite expectations.
\item Expectations exist for all required moments (mean, covariance, etc.).
\item Linear transformations are applied using constant matrices and vectors.
\end{itemize}

\section*{4. Main Results / Theorems}

\textbf{Linear Transformation of Random Vectors}

\[
Y = AX + b
\]

Where:
\begin{itemize}
\item $A$: transformation matrix
\item $b$: constant vector
\item $X$: input random vector
\end{itemize}

Expanded form (for $N$-dimensional case):
\[
Y_i = \sum_{j=1}^{N} a_{ij} X_j + b_i
\]

\textbf{Mean Transformation}

\[
\mu_Y = E[Y]
\]

Final result:
\[
\mu_Y = A\mu_X + b
\]

\textbf{Correlation Matrix Transformation}

\[
R_{YY} = E[YY^T]
\]

Final expanded result:
\[
R_{YY} = AR_{XX}A^T + A\mu_X b^T + b\mu_X^T A^T + bb^T
\]

\textbf{Covariance Transformation (Key Observation)}

\[
Y - \mu_Y = A(X - \mu_X)
\]

This shows that covariance is independent of shift $b$.

\section*{5. Proofs / Derivations}

\textbf{Mean Transformation Derivation}

\[
\mu_Y = E[Y]
\]

Substitute:
\[
Y = AX + b
\]

\[
\mu_Y = E[AX + b]
\]

Using linearity:
\[
\mu_Y = E[AX] + E[b]
\]

Since $A$ is constant:
\[
E[AX] = A E[X] = A\mu_X
\]

\[
E[b] = b
\]

Therefore:
\[
\mu_Y = A\mu_X + b
\]

\textbf{Correlation Matrix Derivation}

\[
R_{YY} = E[YY^T]
\]

Substitute:
\[
Y = AX + b
\]

\[
R_{YY} = E[(AX + b)(AX + b)^T]
\]

Expand:
\[
= E[AXX^T A^T + AXb^T + bX^T A^T + bb^T]
\]

Apply expectation:
\[
R_{YY} = AE[XX^T]A^T + AE[X]b^T + bE[X]^T A^T + bb^T
\]

\[
R_{YY} = AR_{XX}A^T + A\mu_X b^T + b\mu_X^T A^T + bb^T
\]

\textbf{Covariance Insight}

\[
Y - \mu_Y = (AX + b) - (A\mu_X + b)
\]

\[
Y - \mu_Y = A(X - \mu_X)
\]

This proves:
\begin{itemize}
\item Shift does not affect covariance
\item Only linear transformation influences spread
\end{itemize}

\section*{6. Worked Examples}

\textbf{Example: Smart City Traffic + Weather System}

Step 1: Define Random Vector
\[
X = [X_1 \; X_2 \; X_3 \; X_4 \; X_5]
\]

Step 2: Mean Vector
\[
\mu_X = E[X]
\]

Step 3: Covariance Matrix Interpretation

\begin{itemize}
\item More rain $\Rightarrow$ more congestion
\item More rain $\Rightarrow$ lower speed
\item More traffic $\Rightarrow$ higher pollution
\item Higher temperature $\Rightarrow$ pollution disperses
\end{itemize}

Thus, covariance matrix acts as a dependency map of the system.

\section*{7. Summary}

\begin{itemize}
\item Multiple random variables are represented using a random vector.
\item Mean vector captures expected values of each component.
\item Correlation matrix captures second-order relationships.
\item Covariance matrix measures dependencies and variability.
\item Linear transformation: $Y = AX + b$
\item Mean transforms as: $\mu_Y = A\mu_X + b$
\item Correlation transforms with additional terms due to shift
\item Covariance depends only on $A$, not on $b$
\end{itemize}

\end{document}